% Ngân hàng tự sinh (AI)
\begin{bt}%[0H9N1-1]
% nguon: AI soạn tương tự câu mẫu 959cabbe
% mau: 959cabbe
Trên trục tọa độ $\left(O;\overrightarrow{i}\right)$ cho 2 điểm $A$, $B$ có tọa độ lần lượt là $-5; 3$. Tìm tọa độ của véc-tơ $\overrightarrow{AB}$.
	\loigiai{
		Ta có $\overrightarrow{AB}=3+5=8$.\\
		Vậy toạ độ $\overrightarrow{AB}$ là $8$.
	}
\end{bt}

\begin{bt}%[0H9N1-1]
% nguon: AI soạn tương tự câu mẫu 959cabbe
% mau: 959cabbe
Trên trục tọa độ $\left(O;\overrightarrow{i}\right)$ cho 2 điểm $M$, $N$ có tọa độ lần lượt là $-7; -1$. Tìm tọa độ của véc-tơ $\overrightarrow{MN}$.
	\loigiai{
		Ta có $\overrightarrow{MN}=-1+7=6$.\\
		Vậy toạ độ $\overrightarrow{MN}$ là $6$.
	}
\end{bt}

\begin{bt}%[0H9N1-1]
% nguon: AI soạn tương tự câu mẫu 959cabbe
% mau: 959cabbe
Trên trục tọa độ $\left(O;\overrightarrow{i}\right)$ cho 2 điểm $P$, $Q$ có tọa độ lần lượt là $4; 9$. Tìm tọa độ của véc-tơ $\overrightarrow{PQ}$.
	\loigiai{
		Ta có $\overrightarrow{PQ}=9-4=5$.\\
		Vậy toạ độ $\overrightarrow{PQ}$ là $5$.
	}
\end{bt}

\begin{bt}%[0H9H1-1]
% nguon: AI soạn tương tự câu mẫu 959cabbe
% mau: 959cabbe nang_cao
	Trên trục tọa độ $\left(O;\overrightarrow{i}\right)$ cho ba điểm $A$, $B$, $C$ có tọa độ lần lượt là $-3;\ 2;\ 7$. Tìm tọa độ của véc-tơ $\overrightarrow{AB}$ và tọa độ của véc-tơ $\overrightarrow{BC}$, từ đó tính độ dài đoạn thẳng $AC$.
	\loigiai{
		Ta có $\overrightarrow{AB}=2-(-3)=2+3=5$. Vậy tọa độ $\overrightarrow{AB}$ là $5$.\\
		Ta có $\overrightarrow{BC}=7-2=5$. Vậy tọa độ $\overrightarrow{BC}$ là $5$.\\
		Ta có $\overrightarrow{AC}=7-(-3)=7+3=10$.\\
		Vậy độ dài đoạn thẳng $AC=\left|\overrightarrow{AC}\right|=10$.
	}
\end{bt}

\begin{bt}%[0H9N1-2]
% nguon: AI soạn tương tự câu mẫu ae5e3c24
% mau: ae5e3c24
Trong mặt phẳng tọa độ $Oxy$, cho $2$ véc-tơ $\overrightarrow{a}=(5;3)$, $\overrightarrow{b}=(-6;10)$. Tính tích vô hướng $\overrightarrow{a}\cdot \overrightarrow{b}$ và chứng tỏ $\overrightarrow{a}\perp\overrightarrow{b}$.
	\loigiai{
		Ta có $\overrightarrow{a}\cdot \overrightarrow{b} = 5\cdot (-6) + 3\cdot 10 =0$.\\
		Suy ra $\overrightarrow{a}\perp\overrightarrow{b}$.
	}
\end{bt}

\begin{bt}%[0H9N1-2]
% nguon: AI soạn tương tự câu mẫu ae5e3c24
% mau: ae5e3c24
Trong mặt phẳng tọa độ $Oxy$, cho $2$ véc-tơ $\overrightarrow{a}=(7;-2)$, $\overrightarrow{b}=(4;14)$. Tính tích vô hướng $\overrightarrow{a}\cdot \overrightarrow{b}$ và chứng tỏ $\overrightarrow{a}\perp\overrightarrow{b}$.
	\loigiai{
		Ta có $\overrightarrow{a}\cdot \overrightarrow{b} = 7\cdot 4 + (-2)\cdot 14 =0$.\\
		Suy ra $\overrightarrow{a}\perp\overrightarrow{b}$.
	}
\end{bt}

\begin{bt}%[0H9N1-2]
% nguon: AI soạn tương tự câu mẫu ae5e3c24
% mau: ae5e3c24
Trong mặt phẳng tọa độ $Oxy$, cho $2$ véc-tơ $\overrightarrow{a}=(-9;4)$, $\overrightarrow{b}=(8;18)$. Tính tích vô hướng $\overrightarrow{a}\cdot \overrightarrow{b}$ và chứng tỏ $\overrightarrow{a}\perp\overrightarrow{b}$.
	\loigiai{
		Ta có $\overrightarrow{a}\cdot \overrightarrow{b} = (-9)\cdot 8 + 4\cdot 18 =0$.\\
		Suy ra $\overrightarrow{a}\perp\overrightarrow{b}$.
	}
\end{bt}

\begin{bt}%[0H9H1-2]
% nguon: AI soạn tương tự câu mẫu ae5e3c24
% mau: ae5e3c24 nang_cao
	Trong mặt phẳng tọa độ $Oxy$, cho $2$ véc-tơ $\overrightarrow{a}=(8;-6)$, $\overrightarrow{b}=(9;12)$. Tính tích vô hướng $\overrightarrow{a}\cdot \overrightarrow{b}$ và chứng tỏ $\overrightarrow{a}\perp\overrightarrow{b}$.
	\loigiai{
		Ta có $\overrightarrow{a}\cdot \overrightarrow{b} = 8\cdot 9 + (-6)\cdot 12 = 72 - 72 = 0$.\\
		Suy ra $\overrightarrow{a}\perp\overrightarrow{b}$.
	}
\end{bt}

\begin{bt}%[0H9N1-3]
% nguon: AI soạn tương tự câu mẫu 04017960
% mau: 04017960
Trong mặt phẳng tọa độ $Oxy$, cho hai điểm $A(-3; 2)$ và $B(5; 4)$. Tìm tọa độ $\overrightarrow{AB}$.
	\loigiai{
		Ta có $\overrightarrow{AB}=(5-(-3); 4-2)=(8; 2)$.
	}
\end{bt}

\begin{bt}%[0H9N1-3]
% nguon: AI soạn tương tự câu mẫu 04017960
% mau: 04017960
Trong mặt phẳng tọa độ $Oxy$, cho hai điểm $P(2; -5)$ và $Q(6; 1)$. Tìm tọa độ $\overrightarrow{PQ}$.
	\loigiai{
		Ta có $\overrightarrow{PQ}=(6-2; 1-(-5))=(4; 6)$.
	}
\end{bt}

\begin{bt}%[0H9N1-3]
% nguon: AI soạn tương tự câu mẫu 04017960
% mau: 04017960
Trong mặt phẳng tọa độ $Oxy$, cho hai điểm $C(-1; -3)$ và $D(4; 2)$. Tìm tọa độ $\overrightarrow{CD}$.
	\loigiai{
		Ta có $\overrightarrow{CD}=(4-(-1); 2-(-3))=(5; 5)$.
	}
\end{bt}

\begin{bt}%[0H9H1-3]
% nguon: AI soạn tương tự câu mẫu 04017960
% mau: 04017960 nang_cao
Trong mặt phẳng tọa độ $Oxy$, cho ba điểm $A(-2; 3)$, $B(4; -1)$ và điểm $C$ thỏa mãn $\overrightarrow{AC}=(5; 2)$. Tìm tọa độ $\overrightarrow{BC}$.
	\loigiai{
		Gọi $C(x; y)$. Ta có $\overrightarrow{AC}=(x+2; y-3)$.\\
		Vì $\overrightarrow{AC}=(5; 2)$ nên $x+2=5$ và $y-3=2$, suy ra $x=3$ và $y=5$, tức $C(3; 5)$.\\
		Khi đó $\overrightarrow{BC}=(3-4; 5-(-1))=(-1; 6)$.\\
		Vậy $\overrightarrow{BC}=(-1; 6)$.
	}
\end{bt}

\begin{bt}%[0H9H1-4]
% nguon: AI soạn tương tự câu mẫu 2f967af4
% mau: 2f967af4
Trong mặt phẳng $Oxy$, cho $A\left(1; 2\right), B\left(3; 4\right), C\left(5; 0\right)$. Gọi $M$ là điểm thuộc $BC$ sao cho $\overrightarrow{MB}=-2 \overrightarrow{BC}$ và $G$ là trọng tâm của tam giác $ABC$.
\begin{enumerate}
    \item Phân tích $\overrightarrow{AM}$ theo hai véc-tơ $\overrightarrow{AB}$ và $\overrightarrow{AC}$.
    \item Phân tích $\overrightarrow{GM}$ theo hai véc-tơ $\overrightarrow{AB}$ và $\overrightarrow{AC}$.
\end{enumerate}
\loigiai{
\begin{enumerate}
    \item Gọi $M(x;y) \Rightarrow \overrightarrow{MB}=(3-x;4-y); \overrightarrow{BC}=(2;-4); \overrightarrow{AB}=(2;2); \overrightarrow{AC}=(4;-2)$.\\
    Mà $\overrightarrow{MB}=-2 \overrightarrow{BC} \Leftrightarrow \heva{&3-x=-4\\ &4-y=8} \Leftrightarrow \heva{&x=7\\ &y=-4} \Rightarrow M(7;-4) \Rightarrow \overrightarrow{AM}=(6;-6)$.\\
    Giả sử $\overrightarrow{AM} = m\overrightarrow{AB}+n\overrightarrow{AC} \Leftrightarrow \heva{&6=2m+4n\\ &-6=2m-2n} \Leftrightarrow \heva{&m= -1\\&n=2.}$\\
    Vậy $\overrightarrow{AM} = -\overrightarrow{AB}+2\overrightarrow{AC}$.
    \item Ta có tọa độ trọng tâm $G(3;2) \Rightarrow \overrightarrow{GM}=(4;-6)$.\\
    Giả sử $\overrightarrow{GM} = h\overrightarrow{AB}+k\overrightarrow{AC} \Leftrightarrow \heva{&4=2h+4k\\ &-6=2h-2k} \Leftrightarrow \heva{&h= -\dfrac{4}{3}\\ &k= \dfrac{5}{3}.}$\\
    Vậy $\overrightarrow{GM} = -\dfrac{4}{3}\overrightarrow{AB}+\dfrac{5}{3}\overrightarrow{AC}$.
\end{enumerate}
}
\end{bt}

\begin{bt}%[0H9H1-4]
% nguon: AI soạn tương tự câu mẫu 2f967af4
% mau: 2f967af4
Trong mặt phẳng $Oxy$, cho $A\left(2; -1\right), B\left(4; 1\right), C\left(0; 3\right)$. Gọi $M$ là điểm thuộc $BC$ sao cho $\overrightarrow{MB}=-2 \overrightarrow{BC}$ và $G$ là trọng tâm của tam giác $ABC$.
\begin{enumerate}
    \item Phân tích $\overrightarrow{AM}$ theo hai véc-tơ $\overrightarrow{AB}$ và $\overrightarrow{AC}$.
    \item Phân tích $\overrightarrow{GM}$ theo hai véc-tơ $\overrightarrow{AB}$ và $\overrightarrow{AC}$.
\end{enumerate}
\loigiai{
\begin{enumerate}
    \item Gọi $M(x;y) \Rightarrow \overrightarrow{MB}=(4-x;1-y); \overrightarrow{BC}=(-4;2); \overrightarrow{AB}=(2;2); \overrightarrow{AC}=(-2;4)$.\\
    Mà $\overrightarrow{MB}=-2 \overrightarrow{BC} \Leftrightarrow \heva{&4-x=8\\ &1-y=-4} \Leftrightarrow \heva{&x=-4\\ &y=5} \Rightarrow M(-4;5) \Rightarrow \overrightarrow{AM}=(-6;6)$.\\
    Giả sử $\overrightarrow{AM} = m\overrightarrow{AB}+n\overrightarrow{AC} \Leftrightarrow \heva{&-6=2m-2n\\ &6=2m+4n} \Leftrightarrow \heva{&m= -1\\&n=2.}$\\
    Vậy $\overrightarrow{AM} = -\overrightarrow{AB}+2\overrightarrow{AC}$.
    \item Ta có tọa độ trọng tâm $G(2;1) \Rightarrow \overrightarrow{GM}=(-6;4)$.\\
    Giả sử $\overrightarrow{GM} = h\overrightarrow{AB}+k\overrightarrow{AC} \Leftrightarrow \heva{&-6=2h-2k\\ &4=2h+4k} \Leftrightarrow \heva{&h= -\dfrac{4}{3}\\ &k= \dfrac{5}{3}.}$\\
    Vậy $\overrightarrow{GM} = -\dfrac{4}{3}\overrightarrow{AB}+\dfrac{5}{3}\overrightarrow{AC}$.
\end{enumerate}
}
\end{bt}

\begin{bt}%[0H9H1-4]
% nguon: AI soạn tương tự câu mẫu 2f967af4
% mau: 2f967af4
Trong mặt phẳng $Oxy$, cho $A\left(-1; 3\right), B\left(1; 5\right), C\left(3; 1\right)$. Gọi $M$ là điểm thuộc $BC$ sao cho $\overrightarrow{MB}=-2 \overrightarrow{BC}$ và $G$ là trọng tâm của tam giác $ABC$.
\begin{enumerate}
    \item Phân tích $\overrightarrow{AM}$ theo hai véc-tơ $\overrightarrow{AB}$ và $\overrightarrow{AC}$.
    \item Phân tích $\overrightarrow{GM}$ theo hai véc-tơ $\overrightarrow{AB}$ và $\overrightarrow{AC}$.
\end{enumerate}
\loigiai{
\begin{enumerate}
    \item Gọi $M(x;y) \Rightarrow \overrightarrow{MB}=(1-x;5-y); \overrightarrow{BC}=(2;-4); \overrightarrow{AB}=(2;2); \overrightarrow{AC}=(4;-2)$.\\
    Mà $\overrightarrow{MB}=-2 \overrightarrow{BC} \Leftrightarrow \heva{&1-x=-4\\ &5-y=8} \Leftrightarrow \heva{&x=5\\ &y=-3} \Rightarrow M(5;-3) \Rightarrow \overrightarrow{AM}=(6;-6)$.\\
    Giả sử $\overrightarrow{AM} = m\overrightarrow{AB}+n\overrightarrow{AC} \Leftrightarrow \heva{&6=2m+4n\\ &-6=2m-2n} \Leftrightarrow \heva{&m= -1\\&n=2.}$\\
    Vậy $\overrightarrow{AM} = -\overrightarrow{AB}+2\overrightarrow{AC}$.
    \item Ta có tọa độ trọng tâm $G(1;3) \Rightarrow \overrightarrow{GM}=(4;-6)$.\\
    Giả sử $\overrightarrow{GM} = h\overrightarrow{AB}+k\overrightarrow{AC} \Leftrightarrow \heva{&4=2h+4k\\ &-6=2h-2k} \Leftrightarrow \heva{&h= -\dfrac{4}{3}\\ &k= \dfrac{5}{3}.}$\\
    Vậy $\overrightarrow{GM} = -\dfrac{4}{3}\overrightarrow{AB}+\dfrac{5}{3}\overrightarrow{AC}$.
\end{enumerate}
}
\end{bt}

\begin{bt}%[0H9V1-4]
% nguon: AI soạn tương tự câu mẫu 2f967af4
% mau: 2f967af4 nang_cao
Trong mặt phẳng $Oxy$, cho $A\left(2; -3\right), B\left(0; 1\right), C\left(-2; -1\right)$. Gọi $M$ là điểm thuộc đường thẳng $BC$ sao cho $\overrightarrow{MB}=-3\overrightarrow{BC}$ và $G$ là trọng tâm của tam giác $ABC$.
	\begin{enumerate}
		\item Phân tích $\overrightarrow{AM}$ theo hai véc-tơ $\overrightarrow{AB}$ và $\overrightarrow{AC}$.
		\item Phân tích $\overrightarrow{GM}$ theo hai véc-tơ $\overrightarrow{AB}$ và $\overrightarrow{AC}$.
	\end{enumerate}
	\loigiai{
	\begin{enumerate}
		\item Gọi $M(x;y) \Rightarrow \overrightarrow{MB}=(-x;1-y); \overrightarrow{BC}=(-2;-2); \overrightarrow{AB}=(-2;4); \overrightarrow{AC}=(-4;2)$.\\
		Mà $\overrightarrow{MB}=-3\overrightarrow{BC} \Leftrightarrow \heva{&-x=6\\ &1-y=6} \Leftrightarrow \heva{&x=-6\\ &y=-5} \Rightarrow M(-6;-5) \Rightarrow \overrightarrow{AM}=(-8;-2)$.\\
		Giả sử $\overrightarrow{AM} = m\overrightarrow{AB}+n\overrightarrow{AC} \Leftrightarrow \heva{&-8=-2m-4n\\ &-2=4m+2n} \Leftrightarrow \heva{&m=2\\ &n=1.}$\\
		Vậy $\overrightarrow{AM} = 2\overrightarrow{AB}+\overrightarrow{AC}$.
		\item Ta có tọa độ trọng tâm $G(0;-1) \Rightarrow \overrightarrow{GM}=(-6;-4)$.\\
		Giả sử $\overrightarrow{GM} = h\overrightarrow{AB}+k\overrightarrow{AC} \Leftrightarrow \heva{&-6=-2h-4k\\ &-4=4h+2k} \Leftrightarrow \heva{&h=\dfrac{1}{3}\\ &k=\dfrac{4}{3}.}$\\
		Vậy $\overrightarrow{GM} = \dfrac{1}{3}\overrightarrow{AB}+\dfrac{4}{3}\overrightarrow{AC}$.
	\end{enumerate}
	}
\end{bt}

\begin{bt}%[0H9N1-5]
% nguon: AI soạn tương tự câu mẫu 42212c6c
% mau: 42212c6c
	Véc-tơ $\overrightarrow{a}=(3;0)$ được phân tích theo hai véc-tơ đơn vị $\left(\overrightarrow{i};\overrightarrow{j}\right)$ như thế nào?
	\loigiai{
		Ta có $\overrightarrow{a}=(3;0)\Rightarrow\overrightarrow{a}=3\overrightarrow{i}+0\overrightarrow{j}=3\overrightarrow{i}$.\\
	}
\end{bt}

\begin{bt}%[0H9N1-5]
% nguon: AI soạn tương tự câu mẫu 42212c6c
% mau: 42212c6c
	Véc-tơ $\overrightarrow{b}=(0;-5)$ được phân tích theo hai véc-tơ đơn vị $\left(\overrightarrow{i};\overrightarrow{j}\right)$ như thế nào?
	\loigiai{
		Ta có $\overrightarrow{b}=(0;-5)\Rightarrow\overrightarrow{b}=0\overrightarrow{i}+(-5)\overrightarrow{j}=-5\overrightarrow{j}$.\\
	}
\end{bt}

\begin{bt}%[0H9N1-5]
% nguon: AI soạn tương tự câu mẫu 42212c6c
% mau: 42212c6c
	Véc-tơ $\overrightarrow{c}=(2;0)$ được phân tích theo hai véc-tơ đơn vị $\left(\overrightarrow{i};\overrightarrow{j}\right)$ như thế nào?
	\loigiai{
		Ta có $\overrightarrow{c}=(2;0)\Rightarrow\overrightarrow{c}=2\overrightarrow{i}+0\overrightarrow{j}=2\overrightarrow{i}$.\\
	}
\end{bt}

\begin{bt}%[0H9H1-5]
% nguon: AI soạn tương tự câu mẫu 42212c6c
% mau: 42212c6c nang_cao
	Véc-tơ $\overrightarrow{a}=(0;7)$ được phân tích theo hai véc-tơ đơn vị $\left(\overrightarrow{i};\overrightarrow{j}\right)$ như thế nào? Từ đó tính tọa độ của véc-tơ $\overrightarrow{u}=2\overrightarrow{a}+\overrightarrow{i}$.
	\loigiai{
		Ta có $\overrightarrow{a}=(0;7)\Rightarrow\overrightarrow{a}=0\overrightarrow{i}+7\overrightarrow{j}=7\overrightarrow{j}$.\\
		Suy ra $\overrightarrow{u}=2\overrightarrow{a}+\overrightarrow{i}=2\cdot 7\overrightarrow{j}+\overrightarrow{i}=\overrightarrow{i}+14\overrightarrow{j}$.\\
		Vậy $\overrightarrow{u}=(1;14)$.
	}
\end{bt}

\begin{bt}%[0H9H1-6]
% nguon: AI soạn tương tự câu mẫu b279dbb9
% mau: b279dbb9
	Cho $\Delta ABC$ có $A(3; 5)$, $B(2;-3)$, $C(-6; 1)$. Tọa độ điểm $M(a;b)$ thuộc đoạn $BC$ sao cho $S_{\Delta MAB}=3S_{\Delta MAC}$. Tính giá trị của biểu thức $a+b$.
	\loigiai{
		Kẻ $AH \perp BC$. $S_{\Delta MAB}=3S_{\Delta MAC}\Leftrightarrow \dfrac{1}{2}AH\cdot MB=3\cdot \dfrac{1}{2}AH\cdot MC\Leftrightarrow MB=3MC$. \\
		Mà $\overrightarrow{MB}$ và $\overrightarrow{MC}$ ngược hướng $\Rightarrow \overrightarrow{MB}=-3\overrightarrow{MC}$ \\
		$\Leftrightarrow \heva{&x_B-x_M=-3(x_C-x_M)\\&y_B-y_M=-3(y_C-y_M)}\Leftrightarrow\heva{&2-x_M=-3(-6-x_M)\\&-3-y_M=-3(1-y_M)}$ \\
		$\Leftrightarrow\heva{&2-x_M=18+3x_M\\&-3-y_M=-3+3y_M}\Leftrightarrow\heva{&-4x_M=16\\&-4y_M=0}\Leftrightarrow\heva{&x_M=-4\\&y_M=0}$. \\
		Suy ra $M(-4;0)$, do đó $a+b=-4+0=-4$.
	}
\end{bt}

\begin{bt}%[0H9H1-6]
% nguon: AI soạn tương tự câu mẫu b279dbb9
% mau: b279dbb9
	Cho $\Delta ABC$ có $A(7; 2)$, $B(1; 4)$, $C(-3;-2)$. Tọa độ điểm $M(a;b)$ thuộc đoạn $BC$ sao cho $S_{\Delta MAB}=2S_{\Delta MAC}$. Tính giá trị của biểu thức $a+b$.
	\loigiai{
		Kẻ $AH \perp BC$. $S_{\Delta MAB}=2S_{\Delta MAC}\Leftrightarrow \dfrac{1}{2}AH\cdot MB=2\cdot \dfrac{1}{2}AH\cdot MC\Leftrightarrow MB=2MC$. \\
		Mà $\overrightarrow{MB}$ và $\overrightarrow{MC}$ ngược hướng $\Rightarrow \overrightarrow{MB}=-2\overrightarrow{MC}$ \\
		$\Leftrightarrow \heva{&x_B-x_M=-2(x_C-x_M)\\&y_B-y_M=-2(y_C-y_M)}\Leftrightarrow\heva{&1-x_M=-2(-3-x_M)\\&4-y_M=-2(-2-y_M)}$ \\
		$\Leftrightarrow\heva{&1-x_M=6+2x_M\\&4-y_M=4+2y_M}\Leftrightarrow\heva{&-3x_M=5\\&-3y_M=0}\Leftrightarrow\heva{&x_M=-\dfrac{5}{3}\\&y_M=0}$. \\
		Suy ra $M\left(-\dfrac{5}{3};0\right)$, do đó $a+b=-\dfrac{5}{3}+0=-\dfrac{5}{3}$.
	}
\end{bt}

\begin{bt}%[0H9H1-6]
% nguon: AI soạn tương tự câu mẫu b279dbb9
% mau: b279dbb9
	Cho $\Delta ABC$ có $A(2; 8)$, $B(5;-2)$, $C(-7; 4)$. Tọa độ điểm $M(a;b)$ thuộc đoạn $BC$ sao cho $S_{\Delta MAB}=4S_{\Delta MAC}$. Tính giá trị của biểu thức $a+b$.
	\loigiai{
		Kẻ $AH \perp BC$. $S_{\Delta MAB}=4S_{\Delta MAC}\Leftrightarrow \dfrac{1}{2}AH\cdot MB=4\cdot \dfrac{1}{2}AH\cdot MC\Leftrightarrow MB=4MC$. \\
		Mà $\overrightarrow{MB}$ và $\overrightarrow{MC}$ ngược hướng $\Rightarrow \overrightarrow{MB}=-4\overrightarrow{MC}$ \\
		$\Leftrightarrow \heva{&x_B-x_M=-4(x_C-x_M)\\&y_B-y_M=-4(y_C-y_M)}\Leftrightarrow\heva{&5-x_M=-4(-7-x_M)\\&-2-y_M=-4(4-y_M)}$ \\
		$\Leftrightarrow\heva{&5-x_M=28+4x_M\\&-2-y_M=-16+4y_M}\Leftrightarrow\heva{&-5x_M=23\\&-5y_M=-14}\Leftrightarrow\heva{&x_M=-\dfrac{23}{5}\\&y_M=\dfrac{14}{5}}$. \\
		Suy ra $M\left(-\dfrac{23}{5};\dfrac{14}{5}\right)$, do đó $a+b=-\dfrac{23}{5}+\dfrac{14}{5}=-\dfrac{9}{5}$.
	}
\end{bt}

\begin{bt}%[0H9V1-6]
% nguon: AI soạn tương tự câu mẫu b279dbb9
% mau: b279dbb9 nang_cao
Cho $\Delta ABC$ có $A(6; 4)$, $B(2;-5)$, $C(-6; 3)$. Tọa độ điểm $M(a;b)$ thuộc đoạn $BC$ sao cho $S_{\Delta MAB}=3S_{\Delta MAC}$ và $S_{\Delta MBC}=2S_{\Delta MAB}$. Tính giá trị của biểu thức $a+b$.
	\loigiai{
		Kẻ $AH \perp BC$. \\
		Từ $S_{\Delta MAB}=3S_{\Delta MAC}\Leftrightarrow \dfrac{1}{2}AH\cdot MB=3\cdot \dfrac{1}{2}AH\cdot MC\Leftrightarrow MB=3MC$. \\
		Mà $\overrightarrow{MB}$ và $\overrightarrow{MC}$ ngược hướng $\Rightarrow \overrightarrow{MB}=-3\overrightarrow{MC}$ \\
		$\Leftrightarrow \heva{&x_B-x_M=-3(x_C-x_M)\\&y_B-y_M=-3(y_C-y_M)}\Leftrightarrow\heva{&2-x_M=-3(-6-x_M)\\&-5-y_M=-3(3-y_M)}$ \\
		$\Leftrightarrow\heva{&2-x_M=18+3x_M\\&-5-y_M=-9+3y_M}\Leftrightarrow\heva{&4x_M=-16\\&4y_M=4}\Leftrightarrow\heva{&x_M=-4\\&y_M=1}$. \\
		Suy ra $M(-4;1)$, do đó $a+b=-3$. \\
		Thử lại: $MB=\sqrt{(2+4)^2+(-5-1)^2}=\sqrt{72}=6\sqrt{2}$, $MC=\sqrt{(-6+4)^2+(3-1)^2}=\sqrt{8}=2\sqrt{2}$, đúng $MB=3MC$. \\
		Kiểm tra điều kiện $S_{\Delta MBC}=2S_{\Delta MAB}$: $S_{\Delta MBC}=\dfrac{1}{2}AH\cdot BC=\dfrac{1}{2}AH\cdot 4MC=2\cdot\dfrac{1}{2}AH\cdot 2MC=2S_{\Delta MAB}$ (vì $MB=3MC$, $BC=4MC$, $S_{\Delta MAB}=\dfrac{1}{2}AH\cdot 3MC$), thỏa mãn.
	}
\end{bt}

\begin{bt}%[0H9N2-1]
% nguon: AI soạn tương tự câu mẫu 2dc15a56
% mau: 2dc15a56
	Cho hai véc-tơ $\overrightarrow{a}=(2;3),\overrightarrow{b}=(4;1)$.
	\begin{enumerate}
		\item Tìm toạ độ của véc-tơ $2\overrightarrow{a}+3\overrightarrow{b}$.
		\item Tính các tích vô hướng: $\overrightarrow{a} \cdot \overrightarrow{b},(3\overrightarrow{a}) \cdot (2\overrightarrow{b})$.
	\end{enumerate}
	\loigiai{
		\begin{enumerate}
			\item Ta có: $2\overrightarrow{a}+3\overrightarrow{b}=(2 \cdot 2+3 \cdot 4;2 \cdot 3+3 \cdot 1)=(16;9)$.\\
			\item Ta có: $\overrightarrow{a} \cdot \overrightarrow{b}=2 \cdot 4+3 \cdot 1=11$\\
			$3\overrightarrow{a}=(6;9)$ và $2\overrightarrow{b}=(8;2)$ nên $(3\overrightarrow{a}) \cdot (2\overrightarrow{b})=6 \cdot 8+9 \cdot 2=66$.
		\end{enumerate}
	}
\end{bt}

\begin{bt}%[0H9N2-1]
% nguon: AI soạn tương tự câu mẫu 2dc15a56
% mau: 2dc15a56
	Cho hai véc-tơ $\overrightarrow{a}=(3;-1),\overrightarrow{b}=(2;5)$.
	\begin{enumerate}
		\item Tìm toạ độ của véc-tơ $2\overrightarrow{a}+3\overrightarrow{b}$.
		\item Tính các tích vô hướng: $\overrightarrow{a} \cdot \overrightarrow{b},(3\overrightarrow{a}) \cdot (2\overrightarrow{b})$.
	\end{enumerate}
	\loigiai{
		\begin{enumerate}
			\item Ta có: $2\overrightarrow{a}+3\overrightarrow{b}=(2 \cdot 3+3 \cdot 2;2 \cdot (-1)+3 \cdot 5)=(12;13)$.\\
			\item Ta có: $\overrightarrow{a} \cdot \overrightarrow{b}=3 \cdot 2+(-1) \cdot 5=1$\\
			$3\overrightarrow{a}=(9;-3)$ và $2\overrightarrow{b}=(4;10)$ nên $(3\overrightarrow{a}) \cdot (2\overrightarrow{b})=9 \cdot 4+(-3) \cdot 10=6$.
		\end{enumerate}
	}
\end{bt}

\begin{bt}%[0H9N2-1]
% nguon: AI soạn tương tự câu mẫu 2dc15a56
% mau: 2dc15a56
	Cho hai véc-tơ $\overrightarrow{a}=(4;1),\overrightarrow{b}=(-2;3)$.
	\begin{enumerate}
		\item Tìm toạ độ của véc-tơ $2\overrightarrow{a}+3\overrightarrow{b}$.
		\item Tính các tích vô hướng: $\overrightarrow{a} \cdot \overrightarrow{b},(3\overrightarrow{a}) \cdot (2\overrightarrow{b})$.
	\end{enumerate}
	\loigiai{
		\begin{enumerate}
			\item Ta có: $2\overrightarrow{a}+3\overrightarrow{b}=(2 \cdot 4+3 \cdot (-2);2 \cdot 1+3 \cdot 3)=(2;11)$.\\
			\item Ta có: $\overrightarrow{a} \cdot \overrightarrow{b}=4 \cdot (-2)+1 \cdot 3=-5$\\
			$3\overrightarrow{a}=(12;3)$ và $2\overrightarrow{b}=(-4;6)$ nên $(3\overrightarrow{a}) \cdot (2\overrightarrow{b})=12 \cdot (-4)+3 \cdot 6=-30$.
		\end{enumerate}
	}
\end{bt}

\begin{bt}%[0H9H2-1]
% nguon: AI soạn tương tự câu mẫu 2dc15a56
% mau: 2dc15a56 nang_cao
	Cho hai véc-tơ $\overrightarrow{a}=(2;-1),\overrightarrow{b}=(1;4)$.
	\begin{enumerate}
		\item Tìm toạ độ của véc-tơ $3\overrightarrow{a}+2\overrightarrow{b}$.
		\item Tính các tích vô hướng: $\overrightarrow{a} \cdot \overrightarrow{b},\,(2\overrightarrow{a}) \cdot (3\overrightarrow{b})$.
	\end{enumerate}
	\loigiai{
		\begin{enumerate}
			\item Ta có: $3\overrightarrow{a}+2\overrightarrow{b}=(3 \cdot 2+2 \cdot 1;3 \cdot (-1)+2 \cdot 4)=(8;5)$.\\
			\item Ta có: $\overrightarrow{a} \cdot \overrightarrow{b}=2 \cdot 1+(-1) \cdot 4=-2$.\\
			$2\overrightarrow{a}=(4;-2)$ và $3\overrightarrow{b}=(3;12)$ nên $(2\overrightarrow{a}) \cdot (3\overrightarrow{b})=4 \cdot 3+(-2) \cdot 12=-12$.
		\end{enumerate}
	}
\end{bt}

\begin{bt}%[0H9N2-2]
% nguon: AI soạn tương tự câu mẫu 34f6d564
% mau: 34f6d564
	Cho $\overrightarrow{a}=(3; 2)$ và $\overrightarrow{b}=(-1;5)$. Tính $\cos \alpha$, với $\alpha =\left(\overrightarrow{a}, \overrightarrow{b}\right)$.
	\loigiai{
		Ta có $\cos \alpha =\dfrac{\overrightarrow{a}\cdot \overrightarrow{b}}{\left|\overrightarrow{a}\right|\cdot \left|\overrightarrow{b}\right|} =\dfrac{3\cdot (-1)+2\cdot 5}{\sqrt{3^2+2^2}\cdot \sqrt{(-1)^2+5^2}}=\dfrac{7}{\sqrt{338}}$.
	}
\end{bt}

\begin{bt}%[0H9N2-2]
% nguon: AI soạn tương tự câu mẫu 34f6d564
% mau: 34f6d564
	Cho $\overrightarrow{a}=(4; -1)$ và $\overrightarrow{b}=(2;3)$. Tính $\cos \alpha$, với $\alpha =\left(\overrightarrow{a}, \overrightarrow{b}\right)$.
	\loigiai{
		Ta có $\cos \alpha =\dfrac{\overrightarrow{a}\cdot \overrightarrow{b}}{\left|\overrightarrow{a}\right|\cdot \left|\overrightarrow{b}\right|} =\dfrac{4\cdot 2+(-1)\cdot 3}{\sqrt{4^2+(-1)^2}\cdot \sqrt{2^2+3^2}}=\dfrac{5}{\sqrt{221}}$.
	}
\end{bt}

\begin{bt}%[0H9N2-2]
% nguon: AI soạn tương tự câu mẫu 34f6d564
% mau: 34f6d564
	Cho $\overrightarrow{a}=(1; -3)$ và $\overrightarrow{b}=(-2;5)$. Tính $\cos \alpha$, với $\alpha =\left(\overrightarrow{a}, \overrightarrow{b}\right)$.
	\loigiai{
		Ta có $\cos \alpha =\dfrac{\overrightarrow{a}\cdot \overrightarrow{b}}{\left|\overrightarrow{a}\right|\cdot \left|\overrightarrow{b}\right|} =\dfrac{1\cdot (-2)+(-3)\cdot 5}{\sqrt{1^2+(-3)^2}\cdot \sqrt{(-2)^2+5^2}}=\dfrac{-17}{\sqrt{290}}$.
	}
\end{bt}

\begin{bt}%[0H9H2-2]
% nguon: AI soạn tương tự câu mẫu 34f6d564
% mau: 34f6d564 nang_cao
	Cho $\overrightarrow{a}=(3; -2)$ và $\overrightarrow{b}=(1;5)$. Biết $\alpha=\left(\overrightarrow{a},\overrightarrow{b}\right)$, tính $\cos\alpha$ và cho biết $\alpha$ là góc nhọn, vuông hay tù.
	\loigiai{
		Ta có $\cos \alpha =\dfrac{\overrightarrow{a}\cdot \overrightarrow{b}}{\left|\overrightarrow{a}\right|\cdot \left|\overrightarrow{b}\right|} =\dfrac{3\cdot 1+(-2)\cdot 5}{\sqrt{3^2+(-2)^2}\cdot \sqrt{1^2+5^2}}=\dfrac{-7}{\sqrt{13}\cdot\sqrt{26}}=\dfrac{-7}{13\sqrt{2}}$.\\
		Vì $\cos\alpha<0$ nên $\alpha$ là góc tù.
	}
\end{bt}

\begin{bt}%[0H9H2-3]
% nguon: AI soạn tương tự câu mẫu a3548811
% mau: a3548811
	Trong mặt phẳng tọa độ $Oxy$, cho tam giác $ABC$ có $A(2;-1)$, $B(4;3)$ và $C(10;1)$. Tính chu vi tam giác $ABC$.
	\loigiai{
		Ta có $\heva{&\overrightarrow{AB}=(2;4)\\&\overrightarrow{AC}=(8;2)\\&\overrightarrow{BC}=(6;-2)} \Rightarrow \heva{&AB=\sqrt{2^2+4^2}=2\sqrt{5}\\&AC=\sqrt{8^2+2^2}=2\sqrt{17}\\&BC=\sqrt{6^2+(-2)^2}=2\sqrt{10}.}$\\
		Chu vi tam giác $ABC$ là $AB+AC+BC=2\sqrt{5}+2\sqrt{17}+2\sqrt{10}$.
	}
\end{bt}

\begin{bt}%[0H9H2-3]
% nguon: AI soạn tương tự câu mẫu a3548811
% mau: a3548811
	Trong mặt phẳng tọa độ $Oxy$, cho tam giác $ABC$ có $A(-1;2)$, $B(3;5)$ và $C(5;-1)$. Tính chu vi tam giác $ABC$.
	\loigiai{
		Ta có $\heva{&\overrightarrow{AB}=(4;3)\\&\overrightarrow{AC}=(6;-3)\\&\overrightarrow{BC}=(2;-6)} \Rightarrow \heva{&AB=\sqrt{4^2+3^2}=5\\&AC=\sqrt{6^2+(-3)^2}=3\sqrt{5}\\&BC=\sqrt{2^2+(-6)^2}=2\sqrt{10}.}$\\
		Chu vi tam giác $ABC$ là $AB+AC+BC=5+3\sqrt{5}+2\sqrt{10}$.
	}
\end{bt}

\begin{bt}%[0H9H2-3]
% nguon: AI soạn tương tự câu mẫu a3548811
% mau: a3548811
	Trong mặt phẳng tọa độ $Oxy$, cho tam giác $ABC$ có $A(0;1)$, $B(3;-2)$ và $C(9;0)$. Tính chu vi tam giác $ABC$.
	\loigiai{
		Ta có $\heva{&\overrightarrow{AB}=(3;-3)\\&\overrightarrow{AC}=(9;-1)\\&\overrightarrow{BC}=(6;2)} \Rightarrow \heva{&AB=\sqrt{3^2+(-3)^2}=3\sqrt{2}\\&AC=\sqrt{9^2+(-1)^2}=\sqrt{82}\\&BC=\sqrt{6^2+2^2}=2\sqrt{10}.}$\\
		Chu vi tam giác $ABC$ là $AB+AC+BC=3\sqrt{2}+\sqrt{82}+2\sqrt{10}$.
	}
\end{bt}

\begin{bt}%[0H9V2-3]
% nguon: AI soạn tương tự câu mẫu a3548811
% mau: a3548811 nang_cao
	Trong mặt phẳng tọa độ $Oxy$, cho tam giác $ABC$ có $A(2;-3)$, $B(5;1)$ và $C(11;-1)$. Tính chu vi tam giác $ABC$.
	\loigiai{
		Ta có $\heva{&\overrightarrow{AB}=(3;4)\\&\overrightarrow{AC}=(9;2)\\&\overrightarrow{BC}=(6;-2)} \Rightarrow \heva{&AB=\sqrt{3^2+4^2}=5\\&AC=\sqrt{9^2+2^2}=\sqrt{85}\\&BC=\sqrt{6^2+(-2)^2}=2\sqrt{10}.}$\\
		Chu vi tam giác $ABC$ là $AB+AC+BC=5+\sqrt{85}+2\sqrt{10}$.
	}
\end{bt}

\begin{bt}%[0H9N2-4]
% nguon: AI soạn tương tự câu mẫu 0e8eb5bf
% mau: 0e8eb5bf
Trong mặt phẳng tọa độ $Oxy$, cho điểm $A(2;-3)$, đường thẳng $\Delta_1\colon x-2y+7=0$ và đường thẳng $\Delta_2\colon 3x-y-50=0$.
	\begin{enumerate}
		\item Tính góc giữa đường thẳng $\Delta_1$ và $\Delta_2$.
		\item Tính khoảng cách từ điểm $A$ đến đường thẳng $\Delta_2$.
	\end{enumerate}
	\loigiai{
		\begin{enumerate}
			\item $\Delta_1$ có véc-tơ pháp tuyến $\overrightarrow{n}_1=(1;-2)$ và $\Delta_2$ có véc-tơ pháp tuyến $\overrightarrow{n}_2=(3;-1)$.\\
			$\cos(\Delta_1;\Delta_2)=\left|\cos(\overrightarrow{n}_1;\overrightarrow{n}_2)\right|=\dfrac{|1\cdot3+(-2)\cdot(-1)|}{\sqrt{1^2+(-2)^2}\cdot\sqrt{3^2+(-1)^2}}=\dfrac{5}{\sqrt{5}\cdot\sqrt{10}}=\dfrac{\sqrt{2}}{2}$.\\
			Suy ra $(\Delta_1;\Delta_2)=45^\circ$.
			\item $\mathrm{d}(A,\Delta_2)=\dfrac{|3\cdot2-(-3)-50|}{\sqrt{3^2+(-1)^2}}=\dfrac{41}{\sqrt{10}}=\dfrac{41\sqrt{10}}{10}$.
		\end{enumerate}
	}
\end{bt}

\begin{bt}%[0H9N2-4]
% nguon: AI soạn tương tự câu mẫu 0e8eb5bf
% mau: 0e8eb5bf
Trong mặt phẳng tọa độ $Oxy$, cho điểm $A(4;1)$, đường thẳng $\Delta_1\colon 2x+y-9=0$ và đường thẳng $\Delta_2\colon x-3y+20=0$.
	\begin{enumerate}
		\item Tính góc giữa đường thẳng $\Delta_1$ và $\Delta_2$.
		\item Tính khoảng cách từ điểm $A$ đến đường thẳng $\Delta_2$.
	\end{enumerate}
	\loigiai{
		\begin{enumerate}
			\item $\Delta_1$ có véc-tơ pháp tuyến $\overrightarrow{n}_1=(2;1)$ và $\Delta_2$ có véc-tơ pháp tuyến $\overrightarrow{n}_2=(1;-3)$.\\
			$\cos(\Delta_1;\Delta_2)=\left|\cos(\overrightarrow{n}_1;\overrightarrow{n}_2)\right|=\dfrac{|2\cdot1+1\cdot(-3)|}{\sqrt{2^2+1^2}\cdot\sqrt{1^2+(-3)^2}}=\dfrac{1}{\sqrt{5}\cdot\sqrt{10}}=\dfrac{\sqrt{2}}{10}$.\\
			Suy ra $(\Delta_1;\Delta_2)\approx 81{,}9^\circ$.
			\item $\mathrm{d}(A,\Delta_2)=\dfrac{|4-3\cdot1+20|}{\sqrt{1^2+(-3)^2}}=\dfrac{21}{\sqrt{10}}=\dfrac{21\sqrt{10}}{10}$.
		\end{enumerate}
	}
\end{bt}

\begin{bt}%[0H9N2-4]
% nguon: AI soạn tương tự câu mẫu 0e8eb5bf
% mau: 0e8eb5bf
Trong mặt phẳng tọa độ $Oxy$, cho điểm $A(-1;5)$, đường thẳng $\Delta_1\colon 3x-y+6=0$ và đường thẳng $\Delta_2\colon x+2y-40=0$.
	\begin{enumerate}
		\item Tính góc giữa đường thẳng $\Delta_1$ và $\Delta_2$.
		\item Tính khoảng cách từ điểm $A$ đến đường thẳng $\Delta_2$.
	\end{enumerate}
	\loigiai{
		\begin{enumerate}
			\item $\Delta_1$ có véc-tơ pháp tuyến $\overrightarrow{n}_1=(3;-1)$ và $\Delta_2$ có véc-tơ pháp tuyến $\overrightarrow{n}_2=(1;2)$.\\
			$\cos(\Delta_1;\Delta_2)=\left|\cos(\overrightarrow{n}_1;\overrightarrow{n}_2)\right|=\dfrac{|3\cdot1+(-1)\cdot2|}{\sqrt{3^2+(-1)^2}\cdot\sqrt{1^2+2^2}}=\dfrac{1}{\sqrt{10}\cdot\sqrt{5}}=\dfrac{\sqrt{2}}{10}$.\\
			Suy ra $(\Delta_1;\Delta_2)\approx 81{,}9^\circ$.
			\item $\mathrm{d}(A,\Delta_2)=\dfrac{|-1+2\cdot5-40|}{\sqrt{1^2+2^2}}=\dfrac{31}{\sqrt{5}}=\dfrac{31\sqrt{5}}{5}$.
		\end{enumerate}
	}
\end{bt}

\begin{bt}%[0H9H2-4]
% nguon: AI soạn tương tự câu mẫu 0e8eb5bf
% mau: 0e8eb5bf nang_cao
Trong mặt phẳng tọa độ $Oxy$, cho điểm $A(5;-2)$, đường thẳng $\Delta_1\colon 2x-3y+18=0$ và đường thẳng $\Delta_2\colon x+4y-60=0$.
	\begin{enumerate}
		\item Tính góc giữa đường thẳng $\Delta_1$ và $\Delta_2$.
		\item Tính khoảng cách từ điểm $A$ đến đường thẳng $\Delta_2$.
	\end{enumerate}
	\loigiai{
		\begin{enumerate}
			\item $\Delta_1$ có véc-tơ pháp tuyến $\overrightarrow{n}_1=(2;-3)$ và $\Delta_2$ có véc-tơ pháp tuyến $\overrightarrow{n}_2=(1;4)$.\\
			$\cos(\Delta_1;\Delta_2)=\left|\cos(\overrightarrow{n}_1;\overrightarrow{n}_2)\right|=\dfrac{|2\cdot1+(-3)\cdot4|}{\sqrt{2^2+(-3)^2}\cdot\sqrt{1^2+4^2}}=\dfrac{|2-12|}{\sqrt{13}\cdot\sqrt{17}}=\dfrac{10}{\sqrt{221}}$.\\
			Suy ra $(\Delta_1;\Delta_2)\approx 47{,}7^\circ$.
			\item $\mathrm{d}(A,\Delta_2)=\dfrac{|1\cdot5+4\cdot(-2)-60|}{\sqrt{1^2+4^2}}=\dfrac{|5-8-60|}{\sqrt{17}}=\dfrac{63}{\sqrt{17}}=\dfrac{63\sqrt{17}}{17}$.
		\end{enumerate}
	}
\end{bt}
